\par\noindent\textbf{(1)} 对于二元假设密度矩阵ρ0、ρ1与先验各半，定义差算符Δ=(ρ0−ρ1)/2。最佳POVM投影到Δ的正本征子空间判定0，负本征子空间判定1；成功率Psucc=[1+||Δ||\_1]/2，错误率Perr=[1−||Δ||\_1]/2。
\par\noindent\textbf{(2)} 在|ψ0>及其正交补基底中，选|ψ1>=s|ψ0>+√(1−s²)|ψ0\textasciicircum{}⊥>，可计算ρ0−ρ1为迹零2×2 Hermitian矩阵，其本征值为±√(1−s²)，所以Δ特征值±√(1−s²)/2。
\par\noindent\textbf{(3)} 迹范数||Δ||\_1=√(1−s²)，代入得到Perr,min=[1−√(1−s²)]/2。最优测量依赖两态相对方向，并非始终测原始计算基；任意其他二元测量的平均错误率不能低于该值。
\par\noindent\textbf{(4)} s=0两态正交，Perr=0；s=1两态相同，Perr=1/2，猜测与测量没有区分能力。若有n个独立拷贝，两候选总态内积模变为s\textasciicircum{}n，全局最优错误率为[1−√(1−s\textasciicircum{}\{2n\})]/2，随n增大趋零（除s=1）。
